\begin{center}
\begin{tikzpicture}[scale=1.25, line cap=round, line join=round]
  \definecolor{trigblue}{RGB}{30,110,190}
  \coordinate (A) at (0,0);
  \coordinate (B) at (3.4,0);
  \coordinate (C) at (3.4,2.0);
  \draw[very thick] (A)--(B)--(C)--cycle;
  \draw[line width=2.1pt,trigblue] (A)--(B);
  \draw[line width=2.1pt,trigblue] (B)--(C);
  \draw[line width=2.1pt,gray!65] (A)--(C);
  \draw (3.1,0)--(3.1,0.3)--(3.4,0.3);
  \draw[trigblue] (0.55,0) arc[start angle=0,end angle=30.5,radius=0.55];
  \node[trigblue] at (0.78,0.22) {$\theta$};
  \node[below,trigblue] at ($(A)!0.5!(B)$) {$\cos\theta$ side};
  \node[right,trigblue] at ($(B)!0.5!(C)$) {$\sin\theta$ side};
  \node[above left] at ($(A)!0.5!(C)$) {hypotenuse};
\end{tikzpicture}
\end{center}


\begin{definitionbox}{Definition}
\begin{definition}[Right-triangle sine]
\label{def:right-triangle-sine}
Let $T$ be a right triangle and let $\theta$ be one of its acute angles. Write
\[
  \operatorname{opp}(T,\theta),\qquad
  \operatorname{adj}(T,\theta),\qquad
  \operatorname{hyp}(T)
\]
for the side lengths opposite $\theta$, adjacent to $\theta$, and opposite the
right angle, respectively. Define
\[
\sin\theta=\frac{\operatorname{opp}(T,\theta)}{\operatorname{hyp}(T)}.
\]
\end{definition}
\end{definitionbox}
\begin{dependencies}
\begin{itemize}
  \item \hyperref[def:euclid-triangle-angle-classes]{Right-angled triangle}
  \item \hyperref[def:euclid-right-angle-and-perpendicular]{Right angle and perpendicular}
  \item \hyperref[def:euclid-acute-angle]{Acute angle}
\end{itemize}
\end{dependencies}

\begin{definitionbox}{Definition}
\begin{definition}[Acute-angle sine]
\label{def:acute-angle-sine}
For an acute Euclidean angle $\theta$, choose any right triangle $T$ having
$\theta$ as an acute angle and set
\[
  \sin\theta=\frac{\operatorname{opp}(T,\theta)}{\operatorname{hyp}(T)}.
\]
\end{definition}
\end{definitionbox}
\begin{dependencies}
\begin{itemize}
  \item \hyperref[def:right-triangle-sine]{Right-triangle sine}
  \item \hyperref[thm:trig-ratios-similarity-invariant]{Similarity invariance of trigonometric ratios}
\end{itemize}
\end{dependencies}

\begin{center}
\begin{tikzpicture}[scale=1.55, line cap=round, line join=round]
  \definecolor{trigblue}{RGB}{30,110,190}
  \coordinate (O) at (0,0);
  \coordinate (P) at (35:1);
  \coordinate (Q) at ({cos(35)},0);
  \draw[gray!55] (O) circle (1);
  \draw[->] (-1.25,0)--(1.35,0) node[right] {$x$};
  \draw[->] (0,-1.15)--(0,1.25) node[above] {$y$};
  \draw[very thick, trigblue] (Q)--(P);
  \draw[very thick] (O)--(P);
  \draw[dropline] (P)--(Q);
  \draw[trigblue] (0.32,0) arc[start angle=0,end angle=35,radius=0.32];
  \node[trigblue] at (0.32,0.10) {$\theta$};
  \fill (P) circle (0.022) node[above right] {$P=(x_\theta,y_\theta)$};
  \node[right,trigblue] at ($(Q)!0.5!(P)$) {$\sin\theta=y_\theta$};
\end{tikzpicture}
\end{center}


\begin{definitionbox}{Definition}
\begin{definition}[Unit-circle sine]
\label{def:unit-circle-sine}
Let $\theta$ be an angle in standard position, and let $P(\theta)=(x_\theta,
y_\theta)$ be the point where its terminal ray meets the unit circle. Define
\[
  \sin\theta=y_\theta.
\]
\end{definition}
\end{definitionbox}
\begin{dependencies}
\begin{itemize}
  \item \hyperref[def:euclid-circle]{Circle}
  \item \hyperref[def:cartesian-product]{Cartesian Product}
\end{itemize}
\end{dependencies}
